\begin{frame}{$\max$ 뒤에 숨은 ``낙관적 부트스트래핑''}
  $\max_{a'}Q(s',a')$을 쓰는 것은 단순히 ``더 큰 수를 쓰는 것''을
  넘어서, \textbf{아직 한 번도 고르지 않은 행동의 가치까지 갱신에
  끌어들이는} 효과가 있다.
  \vskip0.5em
  \begin{itemize}
    \item 위 손계산에서 $a_2$는 이 스텝에서 실제로 \textbf{고르지
      않}았는데도, Q-learning의 목표값에 $Q(s',a_2)=5$가 들어갔다.
    \item 즉 다음 상태의 ``최선''이라는 \textbf{가정 하나로, 실제로
      겪지 않은 경로의 가치}까지 지금 추정치에 반영.
  \end{itemize}
  \vskip0.5em
  이 ``가정''은 \textbf{두 날을 가진 칼}이다:
  \vskip0.3em
  \begin{itemize}
    \item \kb{장점 --- 탐험을 가속.} 아직 시도하지 않은 행동이 ``좋은 것
      같으면'' 그 $Q(s',a')$가 높게 부트스트랩되어 $\arg\max_a Q(s,a)$가
      곧바로 그 행동을 가리키게 되어 탐험이 자연스레 유도됨.
      SARSA처럼 ``실제로 고른 것''만 갱신하면, 한 번도 시도하지 않은
      행동은 \textbf{영원히 낮은 Q값}에 머물러 탐험이 더딤.
    \item \kb{단점 --- 낙관론이 위험을 과소평가.} ``다음 상태에서 최선은
      $a_2$일 것이다''라는 가정이, 실제로 $\varepsilon$-greedy가 가끔
      ``나쁜'' $a_1$을 고를 수 있다는 \textbf{현실을 무시}.
      CliffWalking에서 Q-learning이 ``절벽 바로 옆'' 최단 경로를 선호하는
      이유가 바로 이 낙관론 --- ``다음 칸에선 최선을 고를 테니 위험한
      칸 옆을 지나도 된다''는 가정이, 학습 도중에는 실제로 절벽에
      떨어지는 사고로 이어진다.
  \end{itemize}
\end{frame}
